\documentclass[spanish,letterpaper,oneside]{article}

\def\documenttitle {¿Cómo hacer visible el desarrollo de habilidades matemáticas?}
\def\documentsubtitle {Temas selectos de matemáticas I · Cohorte SD26}
\def\documentsubject {Actividad 5 (Sesión 12) · Indicadores, Evidencias Observables y Propuesta Didáctica}
\def\documentauthor {Mtro. Martín Jonathan de la Cruz Muñoz}
\def\coursename {Temas selectos de matemáticas I}
\def\coursecode {PFM02}

\def\universityname {Instituto de Investigación, Innovación y Estudios de Posgrado para la Educación del Estado de Nuevo León}
\def\universityfaculty {Maestría en Aprendizaje y Enseñanza de las Matemáticas}
\def\universitydepartment {Temas selectos de matemáticas I}
\def\universitydepartmentimage {img/iiiepe/marca-iiiepe.png}
\def\universitydepartmentimagecfg {height=1.75cm}
\def\universitylocation {Monterrey, Nuevo León}
\def\studentfullname {Mtro. Martín Jonathan de la Cruz Muñoz}
\def\studentid {00874}
\def\studentemail {martin.delacruz@campus.iiiepe.edu.mx}
\def\programperiod {2do periodo, agosto--noviembre de 2026 (SD26)}
\def\programmodality {Modalidad presencial y en línea}
\def\programshort {MAEM}

\def\authortable {
  \begin{tabular}{ll}
    \textbf{Alumno:} & \studentfullname \\
    Matrícula: & \studentid \\
    Correo institucional: & \studentemail \\
    Unidad: & \coursename \\
    Programa: & \programshort \\
    Periodo: & \programperiod \\
    & \\
    \multicolumn{2}{l}{Fecha de realización: 7 de noviembre de 2026} \\
    \multicolumn{2}{l}{\universitylocation}
  \end{tabular}
}

\input{template}

\usepackage{helvet}
\renewcommand{\familydefault}{\sfdefault}
\usepackage{amsmath,amssymb,booktabs,array,longtable,pdflscape,tikz,tabularx,colortbl,pgfplots,enumitem,needspace}
\usepackage[most]{tcolorbox}
\usetikzlibrary{arrows.meta,positioning,calc,fit,backgrounds}
\pgfplotsset{compat=1.18}

\definecolor{iiiepeNavy}{HTML}{163B5C}
\definecolor{iiiepeBlue}{HTML}{1F6F9F}
\definecolor{iiiepeAqua}{HTML}{20B7A8}
\definecolor{iiiepeOrange}{HTML}{F28C28}
\definecolor{iiiepePaper}{HTML}{F6F8F7}
\definecolor{iiiepeInk}{HTML}{1E2A32}
\definecolor{iiiepePale}{HTML}{EEF5F8}
\def\portraittitlecolor {iiiepeNavy}
\setcitestyle{authoryear,open={(},close={)}}

\newcommand{\insertcoverwatermark}{%
  \AddToShipoutPictureBG*{%
    \begin{tikzpicture}[remember picture,overlay]
      \node[opacity=.06,inner sep=0pt] at (current page.center)
      {\includegraphics[width=.58\paperwidth]{img/iiiepe/logo-puntos.png}};
    \end{tikzpicture}%
  }%
}

\newtcolorbox{propuestabox}[1][]{%
  enhanced,breakable,
  colback=iiiepePaper,colframe=iiiepeAqua!80!iiiepeNavy,
  boxrule=1pt,leftrule=4pt,arc=2pt,
  left=8pt,right=8pt,top=6pt,bottom=6pt,
  before skip=6pt,after skip=10pt,
  fonttitle=\bfseries\color{white},coltitle=white,
  colbacktitle=iiiepeNavy,
  attach boxed title to top left={xshift=8pt,yshift=-2pt},
  boxed title style={colback=iiiepeNavy,colframe=iiiepeNavy,arc=1pt},
  title={Propuesta didáctica},#1}

\setlength{\parskip}{5pt}\setlength{\parindent}{0pt}
\begin{document}\insertcoverwatermark\templatePortrait\templatePagecfg\templateFinalcfg\fancyfoot{}\fancyfoot[L]{\scriptsize IIIEPE · MAEM · SD26}\fancyfoot[C]{\scriptsize\thepage}\fancyfoot[R]{\scriptsize Actividad 5}\setstretch{1.25}\raggedbottom

\section{Hacer visible el desarrollo de las ocho habilidades SEP}
La Sesión 12 (presentación interactiva de 49 diapositivas) profundiza en cómo operacionalizar las ocho habilidades matemáticas de la SEP mediante \textbf{indicadores de desempeño} (acciones específicas observables) y \textbf{evidencias observables} (producciones concretas que permiten verificarlas), contrastándolas con retos matemáticos breves en clase.

{\small\setstretch{1.06}\renewcommand{\arraystretch}{1.18}
\begin{longtable}{|>{\raggedright\arraybackslash}p{0.13\linewidth}|>{\raggedright\arraybackslash}p{0.25\linewidth}|>{\raggedright\arraybackslash}p{0.31\linewidth}|>{\raggedright\arraybackslash}p{0.27\linewidth}|}\hline
\rowcolor{iiiepePale} \textbf{Habilidad} & \textbf{Indicador observable} & \textbf{Evidencia observable} & \textbf{Reto de la Sesión 12 resuelto} \\\hline
\endfirsthead\hline
\rowcolor{iiiepePale} \textbf{Habilidad} & \textbf{Indicador observable} & \textbf{Evidencia observable} & \textbf{Reto de la Sesión 12 resuelto} \\\hline
\endhead
\textbf{Calcular} & Selecciona y justifica una estrategia de cálculo. & Presenta el procedimiento explicado y el resultado verificado. & 25\,\% de $360$: (1) $360\div 4 = 90$; (2) 20\,\% ($72$) $+$ 5\,\% ($18$) $= 90$. \\\hline
\textbf{Inferir} & Obtiene conclusiones válidas a partir de datos explícitos e implícitos. & Redacta inferencias respaldadas en elementos de la gráfica o texto. & Gráfica de libros (Lun $4$, Mar $7$, Mié $5$, Jue $10$, Vie $8$): total $34$, promedio $6.8$, pico el jueves. \\\hline
\textbf{Comunicar} & Expresa con claridad ideas, procedimientos y resultados y los justifica. & Explicación ordenada y comprensible con apoyo visual/simbólico. & Mostrar con rectángulo en cuartos por qué $\frac{3}{4} = 0.25+0.25+0.25 = 0.75$. \\\hline
\textbf{Medir} & Determina y compara magnitudes con instrumentos y unidades pertinentes. & Registra dimensiones medidas y calcula la magnitud derivada con unidad. & Medir largo y ancho del pizarrón en $\text{cm}$ o $\text{m}$ y obtener su área en $\text{cm}^2/\text{m}^2$. \\\hline
\textbf{Imaginar} & Representa y transforma mentalmente objetos y movimientos espaciales. & Anticipa y dibuja el resultado de la transformación conservando propiedades. & Rotación horaria $90^\circ$ de la figura ``L'' alrededor del pivote inferior izquierdo. \\\hline
\textbf{Estimar} & Aproxima una medida o resultado mediante referentes conocidos. & Explica el referente razonable utilizado para acotar el valor. & Estimar longitud del segmento sin regla usando ancho de palma o fracción de hoja. \\\hline
\textbf{Generalizar} & Identifica regularidades y formula una regla general aplicable a otros casos. & Reconoce el patrón, escribe la regla simbólica y calcula términos nuevos. & Secuencia de cuadrados $1, 4, 9, 16 \implies a_n = n^2$ ($a_{10} = 100$). \\\hline
\textbf{Deducir} & Obtiene y justifica conclusiones lógicas a partir de propiedades conocidas. & Encadena condiciones dadas, propiedades y conclusión lógica justificada. & Triángulo isósceles ($AB=AC$, $\angle A=40^\circ$): $\angle B=\angle C=(180^\circ-40^\circ)/2=70^\circ$. \\\hline
\end{longtable}}

\section{Propuesta escrita: Habilidad, Contenido, PDA y Evidencia (Base del Trabajo Final)}
Siguiendo los seis apartados obligatorios del \textit{Guión 5} y de las diapositivas 46--47 (\textit{Una habilidad, un PDA, cinco minutos}), se presenta la propuesta didáctica diseñada para su ejecución y grabación en la Sesión 13:

\begin{propuestabox}[title={Propuesta Didáctica Escrita · Sesión 12 (Hacia el Trabajo Final)}]
\subsection*{1. Justificación y vinculación curricular}
\begin{itemize}
  \item \textbf{Nivel educativo y grado:} Bachillerato Tecnológico (CONALEP, semestre inicial / MCCEMS — Saberes y Pensamiento Científico y Pensamiento Matemático).
  \item \textbf{Habilidad matemática SEP seleccionada:} \textbf{Generalizar} (en sinergia con \textit{Inferir} y \textit{Comunicar}): descubrir la regularidad matemática en situaciones experimentales para formular la regla algebraica general $y = k x$ y aplicarla a nuevas condiciones.
  \item \textbf{Contenido matemático:} Variación directamente proporcional $y = k x$ en sistemas hidráulicos ($F_2 = k\cdot F_1$, donde $k = A_2/A_1 = (d_2/d_1)^2$).
  \item \textbf{PDA relacionado:} \textit{Resuelve problemas de variación proporcional directa en contextos de física e ingeniería, representando la relación fuerza--área mediante tablas, expresión algebraica y gráfica lineal en GeoGebra, e interpretando la constante $k$ como factor de amplificación}.
  \item \textbf{Justificación de la relación:} Este PDA exige superar el cálculo aislado de un solo dato para abstraer la relación invariante entre fuerza de entrada $F_1$ y fuerza de salida $F_2$. Al explorar el prototipo físico con distintos diámetros, el alumnado generaliza que la amplificación crece con el \textbf{cuadrado} de la razón de diámetros, vinculando geometría del círculo, proporcionalidad directa y modelación algebraica.
\end{itemize}

\subsection*{2. Indicador de desempeño}
Formula la expresión algebraica de proporcionalidad directa $F_2 = k\cdot F_1$ determinando correctamente $k = (d_2/d_1)^2$, calcula el desplazamiento inverso $\Delta x_2 = \Delta x_1 / k$ y predice nuevos valores en una gráfica lineal que pasa por el origen.

\subsection*{3. Evidencia observable}
Producción conjunta compuesta por: (a) tabla de valores $(F_1, F_2)$ verificada con el prototipo de jeringas, (b) ecuación explícita con cálculo paso a paso de $P = F_1/A_1$ y $k = (d_2/d_1)^2$, (c) gráfica en GeoGebra mostrando la recta $F_2 = 4 F_1$ con el punto del reto $(8\text{ N}, 32\text{ N})$ frente al error lineal $(8\text{ N}, 16\text{ N})$, y (d) respuesta escrita al inciso de generalización cuando $d_2 = 6.0\text{ cm}$ ($k' = 16$).

\subsection*{4. Dificultades o errores frecuentes del alumnado y sus causas}
El error más frecuente consiste en confundir la razón lineal de diámetros $d_2/d_1$ con la razón cuadrática de áreas $(d_2/d_1)^2$. Por ejemplo, con jeringas de $d_1 = 1.5\text{ cm}$, $d_2 = 3.0\text{ cm}$ y $F_1 = 8\text{ N}$, muchos estudiantes asumen intuitivamente que la fuerza solo se duplica ($k = 2 \Rightarrow F_2 = 16\text{ N}$ en lugar de $32\text{ N}$), incurriendo en una subestimación de $-16\text{ N}$ ($50\,\text{\symbol{37}}$ de error). La causa radica en la sobregeneralización de la proporcionalidad lineal simple a magnitudes bidimensionales ($A = \pi d^2/4$); se remedia mostrando explícitamente la cancelación de $\pi/4$ y contrastando ambas rectas en GeoGebra.

\subsection*{5. Estrategia didáctica}
Secuencia ABP articulada con la resolución de problemas de la SEP, las fases de Pólya y técnicas de las \textit{Cien Técnicas Didácticas} de la UnADM sobre representaciones múltiples (física, tabular, algebraica y gráfica):
\begin{enumerate}
  \item \textit{Detonante experimental (Proyectos colaborativos, UnADM \#96 y Heurística de Bruner, UnADM \#69):} Situación real de una excavadora hidráulica construida con jeringas ($d_1 = 1.5\text{ cm}$, $d_2 = 3.0\text{ cm}$, $F_1 = 8\text{ N}$, $\Delta x_1 = 4\text{ cm}$).
  \item \textit{Contraste cognitivo (Autoexplicación, UnADM \#16):} Confrontar las predicciones intuitivas ($k = 2 \Rightarrow 16\text{ N}$) contra el cálculo explícito de áreas circulares para justificar $k = 4 \Rightarrow 32\text{ N}$.
  \item \textit{Generalización y modelación tecnológica (Simulación interactiva, UnADM \#99 y Estudio intercalado, UnADM \#40):} Contrastar $F_2 = 4 F_1$ frente a $F_2 = 2 F_1$ en GeoGebra Classic, verificar $\Delta x_2 = \Delta x_1/k = 1\text{ cm}$ y resolver el reto extendido a $d_2 = 6.0\text{ cm}$ ($k' = 16$).
\end{enumerate}
\end{propuestabox}

\section{Cumplimiento de la Rúbrica de Evaluación (Actividad 5)}
{\small\setstretch{1.08}\renewcommand{\arraystretch}{1.18}
\begin{longtable}{|>{\raggedright\arraybackslash}p{0.27\linewidth}|>{\raggedright\arraybackslash}p{0.69\linewidth}|}\hline
\rowcolor{iiiepePale} \textbf{Criterio de la Rúbrica} & \textbf{Cumplimiento verificado en este documento} \\\hline
\endfirsthead\hline
\rowcolor{iiiepePale} \textbf{Criterio de la Rúbrica} & \textbf{Cumplimiento verificado en este documento} \\\hline
\endhead
Justificación y vinculación curricular & Nivel, grado (Bachillerato CONALEP), contenido, PDA y vínculo con la habilidad \textit{Generalizar} explicados con rigor. \\\hline
Indicador de desempeño & Enunciado específico, observable y evaluable alineado a $F_2 = k\cdot F_1$ y $\Delta x_2 = \Delta x_1/k$. \\\hline
Evidencia observable & Tabla, ecuación con $k=(d_2/d_1)^2$, gráfica GeoGebra $(8,32)$ vs $(8,16)$ y respuesta escrita al inciso (d). \\\hline
Dificultades o errores frecuentes & Identifica el error $d_2/d_1$ vs $(d_2/d_1)^2$ ($16\text{ N}$ vs $32\text{ N}$, $50\,\text{\symbol{37}}$ de subestimación) y explica cómo remediarlo. \\\hline
Estrategia didáctica & Secuencia ABP, fases de Pólya y técnicas UnADM (\#96, \#69, \#16, \#99, \#40) con representaciones múltiples y GeoGebra. \\\hline
Entrega a tiempo & Documento maestro, versión editable Word y presentación compilados y validados. \\\hline
\end{longtable}}

\needspace{5\baselineskip}\section*{Referencias}
{\setstretch{1.15}\setlength{\parskip}{6pt}
\par\hangindent=1.27cm\hangafter=1 Instituto de Investigación, Innovación y Estudios de Posgrado para la Educación del Estado de Nuevo León. (2026). \textit{Sesión 12: ¿Cómo hacer visible el desarrollo de habilidades matemáticas? Presentación interactiva (49 diapositivas), Guión 5 y rúbrica}. EVA Posgrados IIIEPE.
\par\hangindent=1.27cm\hangafter=1 Secretaría de Educación Pública. (1994). \textit{Libro para el maestro. Matemáticas. Educación secundaria}. SEP.
\par\hangindent=1.27cm\hangafter=1 Universidad Abierta y a Distancia de México. (s.f.). \textit{Cien técnicas didácticas para la enseñanza y el aprendizaje en entornos educativos}. UnADM.
}

\end{document}
